\section{Proposal: Atomic Goals and Cumulative Levels}
\label{sec:proposal}

\subsection{System and Adversary Model}
\label{sec:model}

A fusion system consists of $m$ sensors, an aggregator, and an off-system
auditor. Before acquisition, an authenticated and versioned public policy $P$
fixes the device roster, an epoch, the canonical sequence of eligible record
identifiers, and the operator $F$ with its public parameters. It also names the
external anchoring procedure and verification material provisioned before the
run. Policy authentication is a prerequisite: the goals below are relative to
the particular $P$ that the auditor accepts.

Sensors produce the policy-ordered sequence
$O=((d_j,q_j,v_j))_{j=1}^{n}$, where $q_j$ is a device-local sequence number;
the policy need not trust a device wall clock. A manifest $M$ contains the
ordered identifiers and binding commitments to their values. The aggregator
publishes an output $y$, a manifest commitment $c_M$, and evidence $\pi$.
Each run has a unique 16-byte identifier $u$. The auditor's view is
$\mathsf{view}(r)=(P,u,y,c_M,\pi)$ plus the provisioned
verification material. Its judge may also hold private protocol state $\eta$,
including a fresh 32-byte anchor state $a$ and a sampling challenge issued
after observing an anchored $c_M$; $\eta$ is not accepted merely because the
producer copied it into $\pi$.

The anchor signs $H(P)\|u\|a\|c_M$. Its trusted service guarantees fresh $a$
and unique $u$, and the auditor records the anchor receipt before recording
acceptance of $y$; the signature alone is not a timestamp. For sampling, the
auditor then issues a separate unpredictable challenge. A computationally
bounded adversary may control the network
and aggregator, modify stored records, forge an origin label, select which
records reach $M$, or choose $y$. A capture profile separately states whether
sensor application software or isolated secrets may also be compromised.

The fixed schedule covers closed windows. Event-triggered acquisition, churn,
and open-world reports additionally require an authenticated closure rule, such
as signed boundary markers; without one we do not claim policy-complete
coverage. Order-sensitive filters consume values in $P$'s canonical order.

\subsection{Atomic Conditions and Reporting Levels}
\label{sec:levels}

The four atomic conditions for a run are:
\begin{description}
\setlength{\itemsep}{0pt}
\item[G1 (anchored integrity)] an external party bound $u$, fresh auditor state,
and a binding $c_M$ before accepting $y$, so replay or a later entry change is detected;
\item[G2 (credential attribution)] each value commitment in $M$ is
authenticated under the credential assigned by $P$ to the device it names;
\item[G3 (policy-relative coverage)] the ordered identifier sequence in $M$
equals the sequence required by $P$; and
\item[G4 (computation conformance)] $y=F(O_M)$ for policy-ordered openings
$O_M$ of all commitments in $M$.
\end{description}
These conditions are distinct dimensions. The V-levels are their cumulative
prefixes: $\gamma_0$ contains all runs and
$\gamma_i=\bigcap_{j=1}^{i}\mathrm{G}_j$ for $1\leq i\leq4$. This prefix is a
reporting convention chosen because a claim about the policy-required
computation should not silently omit the integrity, attribution, or coverage
obligations on which it relies. It does not imply that a mechanism contributing
G2 or G4 must implement every earlier condition.

We instantiate the goal--judge--tolerance scaffold of Cortier et al.
\cite{cortier2016}. Let $\mathcal{A}_{\rho}$ be the adversaries admitted by a
fixed capture profile $\rho$.

\begin{definition}[Level-$i$ verifiability]
A system $\Sigma$ achieves \textbf{\Vlevel{$i$}} under profile $\rho$ with
tolerance $\delta$ if there is a probabilistic polynomial-time judge $J$ such
that, for every $\mathcal{A}\in\mathcal{A}_{\rho}$,
\[
\Pr[\, r \leftarrow \Sigma^{\mathcal{A}} :
r \notin \gamma_i \wedge J(\mathsf{view}(r),\eta)=\mathsf{accept}\,]\leq\delta,
\]
and $J$ accepts a benign run in $\gamma_i$ with overwhelming probability.
\end{definition}

Table~\ref{tab:levels} names the cumulative prefixes. The qualified names are
intentional: a signature attributes a credential rather than proving an honest
transducer reading, and coverage is complete only relative to $P$. By
construction, $\Vlevel{$i{+}1$}\Rightarrow\Vlevel{$i$}$ under the same profile.
The adjacent counterexamples in Section~\ref{sec:separation} establish that the
added conditions are non-redundant; they do not derive a unique total order from
independence alone.

\begin{table}[t]
\caption{Cumulative run-level reporting convention. Availability, physical
truth, and suitability of $F$ remain outside the hierarchy.}
\label{tab:levels}
\centering
\small
\begin{tabular}{@{}lL{2.8cm}L{2.0cm}L{5.3cm}@{}}
\toprule
& Name & Prefix & What the name does not claim \\
\midrule
\Vzero  & None & no run-level condition & any evidentiary guarantee \\
\Vone   & Anchored & G1 & availability or recoverability of openings \\
\Vtwo   & Credential-attributed & G1--G2 & honest transduction or protected keys \\
\Vthree & Policy-complete & G1--G3 & events outside $P$ or open-window closure \\
\Vfour  & Computation-conformant & G1--G4 & physical truth or that $F$ is appropriate \\
\bottomrule
\end{tabular}
\end{table}

A concrete trace explains the cumulative choice. If $P$ requires sequence
numbers 1 through 60, a manifest containing authentic records 1 through 59
satisfies G2 and fails G3. A proof that $y=F(O_M)$ over that incomplete manifest
could satisfy G4 in isolation but is not cumulative V4. Conversely, a
homomorphic authenticator may contribute G2+G4 without constituting a complete
V4 stack.

\subsection{Boundaries and Exclusions}
\label{sec:rationale}

Mechanisms do not name levels. Defining V1 as ``a Merkle tree,'' for example,
would make cost analysis circular and exclude equivalent constructions. The
ALCOA+ vocabulary offers a regulatory reading, not a derivation: its
Attributable, Complete, Enduring, and Accurate terms resemble parts of G1--G4,
but legibility and availability are operational, consistency reaches across
runs, and contemporaneity needs trusted time \cite{mhra2018}. We stop at G4
because $P$ already fixes $F$ and its parameters. Software-build provenance is
an adjacent object, as SLSA illustrates \cite{slsa}.

Capture also bounds what evidence can establish. A sensor controlling its own
measurement path can commit and authenticate a plausible but false value, which
later processing evidence cannot distinguish from an honest one
\cite{przydatek2003sia}. Availability is separate too: a
commitment proves binding only when an opening is supplied; it does not recover
a discarded record.

\subsection{Capture-Resilience Profiles}
\label{sec:rot}

Physical capture changes which parties and secrets the judge may trust, not the
run-level condition. Extracting a sensor's signing key undermines G2 for that
sensor, while compromising only the aggregator need not do so when sensors sign
their own commitments. Every result must therefore pair its V-level with a
capture profile.

Table~\ref{tab:profiles} gives templates a deployment must refine.
A $\rho_{\mathrm{eval}}$ claim must name the evaluated duration and attacks;
hardware supports it only within that scope. RATS appraises evidence against
references, endorsements, and policy \cite{rfc9334}; APEX addresses software
compromise, not general physical resistance \cite{apex2020}.

\begin{table}[t]
\caption{Capture-resilience templates reported independently of a V-level.}
\label{tab:profiles}
\centering
\small
\begin{tabular}{@{}lL{4.3cm}L{4.2cm}@{}}
\toprule
Profile & Adversary admitted & Evidence base required \\
\midrule
$\rho_{\mathrm{soft}}$ & network and aggregator control; no enrolled sensor
application compromise & sensor credential in ordinary application memory \\
$\rho_{\mathrm{iso}}$ & full application-software compromise & isolated key,
freshness, and measured acquisition/execution path \\
$\rho_{\mathrm{eval}}$ & named physical attacks within an evaluation scope &
laboratory-evaluated root of trust and key protection \\
\bottomrule
\end{tabular}
\end{table}

No profile makes a physical observation true. A protected path may stop
compromised software from changing a reading after acquisition, but cannot stop
environmental spoofing or a compromised transducer from supplying a legal-
looking value. Cross-sensor trust estimation addresses this residual risk
\cite{mate2025}; it composes with the evidence goals but is not implied by them.
